\label{app:consentimiento}

Por medio de la firma del presente consentimiento informado \textbf{autoriza de manera completamente voluntaria} el uso de toda la información proporcionada durante esta entrevista y encuestas de seguimiento, así como de su información de contacto, en el desarrollo del trabajo de tesis titulado \textbf{Rediseño de los manipuladores de un robot de servicio de propósito general} elaborada por Rafael Cu\'ellar Ram\'irez y Marco Elian Soriano Pimentel, quienes fungen como investigadores principales y serán las únicas dos personas con acceso a la información recabada. El trabajo se desarrolla en el Laboratorio de Bio-Robótica adscrito al Departamento de Procesamiento de Señales de la División de Ingeniería Eléctrica de la Facultad de Ingeniería de la UNAM, como único centro de investigación participante.

El trabajo de tesis tiene como objetivo el diseño de los manipuladores y efectores finales de un robot de servicio de propósito general, así como el algoritmo de control, para permitirle tomar objetos de uso doméstico.

Como parte de la metodología de diseño centrada en el usuario de Karl T. Ulrich y Steven D. Eppingerm, se requiere conocer la problemática que se pretende resolver a través de las experiencias de los usuarios potenciales y/o de aquellos que ya hayan utilizado tecnologías similares. Con este fin se realizará una entrevista inicial con una duración aproximada de una hora y encuestas de seguimiento posteriores. Durante la entrevista se grabará audio y se tomarán notas. Las encuestas de seguimiento se enviarán por medios electrónicos.

Su participación en este trabajo de investigación no representa riesgo de lesiones o problemas de salud relacionados, por lo tanto, no se contemplan normas de indemnización. Sin embargo, su participación conlleva los riesgos previsibles y no previsibles que se listan a continuación: 

\begin{itemize}
\item Posible incomodidad al compartir información sobre necesidades personales de asistencia doméstica.
\item Fatiga durante la sesión de entrevista.
\item Riesgo mínimo de identificación indirecta a pesar de la codificación de datos.
\end{itemize}

Por otro lado, aunque \textbf{no se contempla una remuneración económica}, su participación conlleva los siguientes beneficios:

\begin{itemize}
\item Contribución directa al desarrollo de tecnología de asistencia robótica que puede mejorar la calidad de vida de personas con necesidades similares.
\item Oportunidad de reflexionar sobre sus experiencias y necesidades, lo cual puede generar aprendizajes valiosos para su propio contexto.
\item Retroalimentación que puede enriquecer su perspectiva sobre el diseño centrado en el usuario.
\item Participación en investigación que busca cerrar la brecha entre el desarrollo tecnológico y las necesidades reales de los usuarios finales.
\end{itemize}

%% Condiciones de confidencialidad y gestión de la información
Se garantizará la privacidad de los participantes mediante consentimiento informado, y la confidencialidad de la información será protegida a través de la codificación de datos, almacenamiento seguro y acceso restringido únicamente al equipo investigador. La informaci\'on ser\'a utilizada con fines estrictamente acad\'emicos y de investigaci\'on y, una vez concluido el trabajo, todas las grabaciones serán eliminadas.

Al autorizar su participación, el equipo investigador se compromete a compartir con usted los resultados a través de los trabajos derivados que sean publicados.

En caso de que tenga alguna duda o inquietud relacionada con su participación, puede contactar a los investigadores principales. Adicionalmente, \textbf{puede retirarse libremente en cualquier momento sin consecuencia alguna}, mientras la investigación esté en curso mediante la \textbf{Carta de revocación de consentimiento} adjunta.

El equipo de investigación se reserva el derecho de terminar su participación sin su consentimiento en caso de que la información proporcionada no resulte relevante para el desarrollo del tema de investigación.

Este proyecto está avalado por el \textbf{Comité de Ética en Investigación y Docencia de la Facultad de Ingeniería}.